\chapter{Literature Review}\label{ch:literature}

\section{Laptop Selection as a Multi-Criteria Decision Problem}\label{sec:mcdm-problem}

Choosing a laptop requires weighing attributes that cannot be reduced to a single scale.
Processing performance, battery life, durability, portability, and price move in different
directions and often trade off against one another, since a lighter chassis tends to shrink
battery capacity and a lower price tends to constrain processing power. \textcite{sonmez_cakir_determination_2020}
frame this directly, noting that new features are added to laptops every year and that the
resulting variety of brands, models, and configurations forces buyers into a comparison problem
that intuition alone cannot resolve consistently. \textcite{gaur_application_2023} reach a similar
conclusion from the demand side, arguing that as shipments grew sharply during the pandemic, the
number of criteria buyers had to reconcile, from brand and price to processor and screen size, grew with it.
A decision problem of this kind, where several conflicting criteria must be evaluated jointly
against a set of alternatives, falls within the scope of Multi-Criteria Decision-Making (MCDM),
the methodological family this chapter now turns to.

\section{An Overview of MCDM Methods}\label{sec:mcdm-overview}

MCDM research has grown into an established branch of operations research. Based on a Scopus
bibliometric analysis, \textcite{zyoud_bibliometric-based_2017} report 10{,}188 documents on AHP
and 2{,}412 documents on TOPSIS through 2015, with China producing the largest national share of
each (34.5\% and 35.1\%, respectively) and \emph{Expert Systems with Applications} standing out as
the most productive outlet for both techniques. This scale of activity reflects a broad family of
methods rather than a single dominant technique.

As an organising lens, this study adopts the five-group taxonomy used by
\textcite{villalba_review_2024} in their domain review and reads it alongside the broader
cross-domain evidence assembled by \textcite{basilio_systematic_2022}. The five families,
laid out in \Cref{fig:mcdm-taxonomy}, differ less in the problems they address than in the
arithmetic they use to reach a ranking. At one end, direct scoring methods such as Simple
Additive Weighting and its extension COPRAS aggregate weighted performance scores into a
single index, and the utility and value family, including MAUT, MAVT, and MIVES, does much
the same after first routing each criterion through an explicit utility function.
Distance-based methods, principally TOPSIS and VIKOR, take a geometric view instead,
ranking alternatives by how close they lie to an ideal solution, whereas pairwise
comparison methods led by AHP, alongside the Analytic Network Process and MACBETH, work
from relative judgements between pairs rather than from absolute scores at all. Outranking
methods such as PROMETHEE and ELECTRE occupy a different position again, building a
preference relation from the degree to which one alternative outranks another rather than
collapsing performance into a single number. This taxonomy is useful because it shows that
AHP and TOPSIS belong to different groups, pairwise comparison and distance-based
respectively, so that combining them addresses two distinct parts of the decision problem,
weighting and ranking, rather than duplicating the same mechanism twice.

\begin{figure}[htbp]
  \centering
  \begin{tikzpicture}[
    font=\scriptsize,
    root/.style={draw, rounded corners, align=center, minimum height=7mm,
                 font=\footnotesize\bfseries, fill=gray!12, inner sep=4pt},
    fam/.style={draw, rounded corners, align=center, text width=2.35cm,
                minimum height=15mm, inner sep=3pt},
    hl/.style={fam, very thick, fill=gray!18},
    lk/.style={-{Latex[length=1.6mm]}, gray!55}
  ]
    \node[fam] (score) at (0,0)    {\textbf{Direct scoring}\\[1pt] SAW, COPRAS};
    \node[fam] (util)  at (2.9,0)  {\textbf{Utility / value}\\[1pt] MAUT, MAVT, MIVES};
    \node[hl]  (dist)  at (5.8,0)  {\textbf{Distance-based}\\[1pt] \textbf{TOPSIS}, VIKOR};
    \node[hl]  (pair)  at (8.7,0)  {\textbf{Pairwise comparison}\\[1pt] \textbf{AHP}, ANP, MACBETH};
    \node[fam] (out)   at (11.6,0) {\textbf{Outranking}\\[1pt] PROMETHEE, ELECTRE};
    \node[root] (root) at (5.8,2.1) {MCDM methods};
    \foreach \n in {score,util,dist,pair,out}{\draw[lk] (root) -- (\n.north);}
    \draw[-{Latex[length=1.8mm]}, thick]
      (pair.south) .. controls +(0,-9mm) and +(0,-9mm) ..
      node[below=1pt, font=\scriptsize\itshape, text width=3.6cm, align=center]
      {this study: AHP weights feed TOPSIS ranking} (dist.south);
  \end{tikzpicture}
  \caption{The five families of MCDM methods, with the pairwise-comparison and
  distance-based groups combined in this study.}
  \label{fig:mcdm-taxonomy}
  \floatnote{\emph{Note:} representative methods are listed under each family; AHP and
  TOPSIS, used in this study, are shown in bold.}{\emph{Source:} adapted from
  \textcite{villalba_review_2024}.}
\end{figure}

The prevalence of these two techniques is confirmed by counts of published applications.
\textcite{basilio_systematic_2022}, reviewing 23{,}494 documents published between 1977 and 2022,
record 6{,}835 applications of AHP and 4{,}907 of TOPSIS, ahead of VIKOR (1{,}475), PROMETHEE
(1{,}382), ANP (1{,}262), ELECTRE (1{,}005), and DEMATEL (888).
\textcite{villalba_review_2024} report a comparable pattern within their own building-assessment
domain, where the distance-based and pairwise comparison groups identified above are the two most
frequently applied, AHP alone accounts for 35\% of the reviewed articles and its combination with
TOPSIS for a further 31\%, ahead of the next most common technique, Simple Additive Weighting, at
11\%, and AHP overall, whether used alone or hybridised with another technique, appears in 71\% of
the sample.
Hybridisation with another MCDM method occurs in 55 of the reviewed articles, while only 5 propose
a modification to an existing method rather than applying it as originally formulated, indicating
that researchers in this domain favour combining established techniques over developing new ones.
\textcite{tighnavard_balasbaneh_systematic_2025} likewise observe that combining AHP and TOPSIS is
used more often than other method pairings, although they caution that no single combination has
emerged as a settled consensus across application domains. Taken together, these figures indicate
that AHP and TOPSIS are not only individually common but are also the pairing most frequently
chosen when a study needs both criteria weights and an alternative ranking, which motivates a
closer look at each method in turn.

\section{The Analytic Hierarchy Process}\label{sec:ahp-theory}

The Analytic Hierarchy Process was introduced by \textcite{saaty_analytic_1987} as a method of
measurement built on ratio scales derived from pairwise comparison. According to \textcite{zyoud_bibliometric-based_2017},
AHP rests on three principles. Decomposition structures a decision into a hierarchy of a
goal, criteria, and alternatives; comparative judgement elicits the relative importance of
elements at the same hierarchical level through pairwise comparison; and synthesis of priorities
aggregates these judgements into an overall priority vector. Comparisons are expressed on a
nine-point ratio scale (\Cref{tab:saaty}), which allows both quantitative and qualitative criteria
to be evaluated within the same framework, a property that makes AHP well suited to decisions such
as laptop selection in which some attributes, for instance processor performance, are naturally
numeric while others, such as brand perception, are not.

Because the judgements underlying a pairwise comparison matrix are elicited from a human
respondent rather than measured directly, AHP incorporates an explicit check on their internal
coherence. The consistency ratio, obtained from the principal eigenvalue of the comparison matrix
as detailed in \Cref{sec:stage2} (\Cref{eq:cr}) and formalised by \textcite{Saaty2008DECISIONMW},
flags matrices in which the respondent's judgements
contradict one another, for example by rating criterion A more important than B, B more important
than C, yet C more important than A. This consistency check is one of the reasons AHP remains
attractive for eliciting subjective expert weights, because unlike a simple ranking or scoring
exercise it provides a quantitative signal for when a judgement set should be revisited. The method is not
without limitations, however. \textcite{zyoud_bibliometric-based_2017} note that AHP is susceptible
to rank reversal, whereby the relative ranking of two existing alternatives can change when a new,
unrelated alternative is added to the comparison, and that the pairwise comparison procedure
assumes criteria are independent of one another, an assumption that does not always hold in
practice.

Where expert judgements are too imprecise to be expressed as single crisp numbers, several studies
extend AHP into a fuzzy environment. \textcite{liu_review_2020} organise this literature around
four technical choices that a fuzzy AHP model must make. It must decide how the relative
importance of a comparison is represented, typically as a fuzzy number rather than a crisp value;
how judgements from multiple experts or multiple criteria are aggregated; how the resulting fuzzy
priorities are defuzzified back into crisp weights for practical use; and how consistency is
measured once judgements are no longer single numbers. This layer of fuzzification is warranted when the study
is designed to preserve linguistic uncertainty rather than require a single comparison value. The
present study instead asks one representative expert to provide explicit Saaty-scale judgements
and reports the resulting matrix in full. Crisp AHP is therefore selected for parsimony and
auditability, not because uncertainty is assumed to be absent. This choice also bounds the
interpretation, since the priority vector represents one internally consistent expert judgement
and should not be read as a group consensus.

The independence assumption and the possibility of rank reversal remain relevant design
limitations. The model treats the retained criteria as analytically separable even where practical
relationships exist, such as those between battery capacity and weight. AHP is used only to weight
criteria rather than to rank the laptop alternatives, which reduces but does not eliminate the
need to examine how the final ranking responds to the chosen criteria structure and weights.

\section{The TOPSIS Method}\label{sec:topsis-theory}

The Technique for Order of Preference by Similarity to Ideal Solution was proposed by
\textcite{Hwang1981MultipleAD} on the principle that the best alternative should lie
simultaneously closest to a positive ideal solution and farthest from a negative ideal solution,
where the positive ideal solution combines the best attainable value on every criterion and the
negative ideal solution combines the worst. \textcite{behzadian_state--art_2012} note that this
requires attribute values that are numeric, monotonically increasing or decreasing in desirability,
and measured on commensurable scales, conditions that a laptop dataset of prices, benchmark scores,
and physical measurements satisfies directly.

Reviewing 266 applications of TOPSIS published across 103 journals since 2000,
\textcite{behzadian_state--art_2012} find the method concentrated in supply chain management and
logistics (27.5\% of the sample) and in design, engineering, and manufacturing systems (23\%),
together accounting for more than half of the applications surveyed. The same review attributes
this spread to several practical advantages. TOPSIS makes full use of the available attribute
information rather than discarding it at an intermediate step, produces a cardinal ranking of every
alternative rather than a partial order, does not require the criteria to be preferentially
independent, and remains conceptually and computationally simple regardless of how many criteria
or alternatives are involved. \textcite{zyoud_bibliometric-based_2017} similarly place TOPSIS as
the second most frequently used MCDM technique after AHP, consistent with the frequency counts
reported in \Cref{sec:mcdm-overview}.

These advantages do not make a TOPSIS ranking invariant. The method is compensatory, in that
strong performance on one criterion can offset weak performance on another once both enter the
weighted distance calculation. Its result also depends on the normalisation rule, the weight
vector, and the observed alternative set from which the ideal and anti-ideal values are
constructed. A change to the candidate set can therefore alter the reference points and the
relative closeness of existing alternatives. In this study, the response is procedural. The
candidate set and collection date are fixed, one normalisation rule is used throughout, the raw
and transformed decision matrices are disclosed, and small differences among the leading closeness
coefficients are interpreted as a short list rather than as evidence of an absolute winner.

As with AHP, a fuzzy extension exists for settings in which attribute values themselves are
imprecise rather than directly measurable. \textcite{chen_extensions_2000} extend TOPSIS to a fuzzy
environment by representing ratings and weights as triangular fuzzy numbers, computing distances to
a fuzzy positive and negative ideal solution through the vertex method, and deriving a closeness
coefficient from these distances. This extension is most useful when performance data are
themselves elicited as linguistic judgements. The laptop dataset used in this study, in contrast,
is built from directly observable specifications and market prices, so the crisp TOPSIS procedure
described in \Cref{sec:stage3} is applied without a fuzzy layer.

\section{Prior Studies on Laptop and Electronics Selection}\label{sec:prior-studies}

Applying MCDM to laptop and related electronics choices is an established but fragmented line of
research, summarised in \Cref{tab:prior-studies}.

\begin{table}[H]
  \centering
  \caption{Prior studies on laptop and electronics selection using MCDM.}
  \label{tab:prior-studies}
  \footnotesize
  \setlength{\tabcolsep}{4pt}
  \begin{tabularx}{\textwidth}{@{}L{3.1cm}L{3.4cm}L{3.6cm}Y@{}}
    \toprule
    \textbf{Source} & \textbf{Context / data} & \textbf{Criteria input} & \textbf{Weighting / ranking} \\
    \midrule
    \textcite{sonmez_cakir_determination_2020} & Laptop; 23 experts & Literature + expert refinement & Fuzzy AHP; DEMATEL; no product ranking \\
    \textcite{gaur_application_2023} & Nine laptop brands; pandemic setting & Nine technical and market criteria & AHP + BWM weights; TOPSIS ranking \\
    \textcite{weng_siew_lam_evaluation_2023} & Mobile-phone alternatives & Technical, physical, and user-related criteria & AHP weights; TOPSIS ranking \\
    \textcite{maghsoudi_hybrid_2026} & Notebook reviews and expert knowledge & Review-derived product features & Sentiment + fuzzy weighting + TOPSIS \\
    \textcite{rau_optimal_2018} & Notebook purchase timing & Value, brand, cost, and eco-efficiency & Scenario weights; TOPSIS ranking \\
    \textcite{liao_determinants_2022} & Eco-friendly laptop experiment & Price, battery, shell, CPU, and display & Conjoint analysis \\
    \textcite{kang_study_2022} & Online laptop reviews & Product, logistics, and after-sales attributes & Empirical review mining \\
    \textcite{guan_cognitive_2022} & Personalised laptop recommendation & Product attributes and user preferences & Pairwise rating + clustering \\
    \textcite{pramanik_comparative_2021} & Mobile computing resources & Computing-resource parameters & Comparative MCDM ranking \\
    \textcite{sostar_assessment_2023} & General consumer behaviour & Cultural, social, personal, psychological factors & AHP weighting \\
    \textcite{jimenez-parra_key_2014} & Remanufactured laptops in Spain & Attitude, norms, and motivation & Behavioural model \\
    \bottomrule
  \end{tabularx}
  \floatsource{\emph{Source:} Authors.}
\end{table}

The evidence in \Cref{tab:prior-studies} shows agreement on the multidimensional character of the
decision but disagreement on which dimension should dominate. \textcite{sonmez_cakir_determination_2020}
place brand image first and show that brand, price, and peripheral properties account for roughly
two-thirds of total priority. By contrast, \textcite{liao_determinants_2022} identify price as the
most influential attribute in an eco-friendly laptop experiment, while
\textcite{weng_siew_lam_evaluation_2023} assign the greatest technical importance to battery life
in their mobile-phone model. These differences are not measurement noise alone; they reflect
different products, populations, criteria definitions, and elicitation procedures.

The studies also confirm that integrated ranking is not methodologically novel by itself.
\textcite{gaur_application_2023} combine AHP and BWM weights with TOPSIS to rank laptop brands, and
\textcite{rau_optimal_2018} use TOPSIS with real notebook data to rank purchase options over time.
Related work extends the evidence base beyond MCDM. \textcite{kang_study_2022} show that online
purchase decisions incorporate logistics and after-sales experience alongside product
configuration, while \textcite{jimenez-parra_key_2014} demonstrate that attitudes and social norms
matter in the remanufactured-laptop context. A technically complete ranking can therefore remain
contextually incomplete if its criteria are fixed before the intended users and local procurement
conditions are examined.

\section{Research Gap and Positioning}\label{sec:gap}

The gap is consequently located in evidence integration rather than in the invention of a new
hybrid method. Existing laptop-ranking studies demonstrate that AHP, BWM, fuzzy weighting, and
TOPSIS can be combined, but they generally begin with criteria selected for a different user group
or study context. Consumer studies identify additional economic, experiential, and service
attributes but do not carry them through to an auditable alternative ranking. The present study
links these streams by separating four decisions that prior work often compresses, namely which
criteria enter the candidate set, which survive screening by the target population, how their
relative weights are elicited, and how a fixed alternative snapshot is ranked.

This positioning also defines what the study does not claim. It applies an established AHP--TOPSIS
pairing rather than proposing a new algorithm, and it evaluates internationally recommended models
priced through Amazon.com rather than claiming coverage of the Vietnamese retail market. Its
contribution lies in the traceable calibration of the criteria to Vietnamese office work and in
the disclosure of the evidence and calculations connecting that calibration to the final ranking.
